\section{Starting electrothermal model}
\label{sec:starting_model}

The starting point of our simulations combines two inherited model components.
Zhang \emph{et al.} formulated the electrical ordinary differential equation
and lumped thermal stochastic differential equation for a \VOtwo{} neuristor
\cite{zhang2024collective}, whereas Almeida \emph{et al.} developed the
history-dependent resistance law used to represent its thermal hysteresis
\cite{almeida2002hysteretic}. We briefly state these equations for a single
neuristor, omitting device indices, before deriving our current-source
formulation.

\begin{figure}[H]
    \centering
    \begin{circuitikz}
        \draw (0,0) node[ocirc]{} node[left]{\(V^{\mathrm{in}}\)}
              to[R, l=\(R^{\mathrm{load}}\)] (3,0)
              node[circ]{} node[above left]{\(V\)} coordinate (Vnode);

        \coordinate (LeftTop) at (3.0,0.9);
        \coordinate (LeftBot) at (3.0,-0.9);
        \coordinate (RightTop) at (5.4,0.9);
        \coordinate (RightMid) at (5.4,0.0);
        \coordinate (RightBot) at (5.4,-0.9);

        \draw (LeftTop) -- (Vnode) -- (LeftBot);
        \draw (RightTop) -- (RightMid) -- (RightBot);
        \draw (LeftTop) to[C, l=\(C\)] (RightTop);
        \draw (LeftBot) to[R, l={\(R(T,\mathcal{H})\)}] (RightBot);
        \draw (RightMid) -- (6.2,0) -- (6.2,-2.2) node[ground]{};
    \end{circuitikz}
    \caption{Single-neuristor equivalent circuit inherited from Zhang
    \emph{et al.} The \VOtwo{} element is represented by capacitive and
    hysteretic resistive branches connected in parallel.}
    \label{fig:single_neuristor_model}
\end{figure}

\subsection{Electrical balance}

The original circuit is voltage driven. A source $V^{\mathrm{in}}(t)$ is
connected to the neuristor node through a load resistance
$R^{\mathrm{load}}$. The \VOtwo{} element is represented by a capacitance
$C$ in parallel with its temperature- and history-dependent resistive branch
$R$. Applying Kirchhoff's current law at the node gives

\begin{equation}
    C\frac{\dd V}{\dd t}
    =
    \frac{V^{\mathrm{in}}-V}{R^{\mathrm{load}}}
    -\frac{V}{R}.
    \label{eq:voltage_balance_physical}
\end{equation}

The first term in Eq.~\eqref{eq:voltage_balance_physical} is the current
delivered by the source through the load resistor. The second is the current
through the resistive branch of the \VOtwo{} element. Their difference charges
or discharges its capacitive branch.

\subsection{Thermal balance}

The neuristor temperature follows the lumped energy balance

\begin{equation}
    C_{\mathrm{th}}\frac{\dd T}{\dd t}
    =
    \frac{V^2}{R}
    -S_e\left(T-T_0\right)
    +\sigma\eta(t).
    \label{eq:thermal_balance}
\end{equation}

Here $C_{\mathrm{th}}$ is the effective thermal capacitance, $V^2/R$ is the
Joule-heating power, and $S_e(T-T_0)$ is the heat transferred from the device
to its surroundings at temperature $T_0$. The term $\sigma\eta(t)$ represents
optional thermal noise; we set $\sigma=0$ in the deterministic simulations.

\subsection{Hysteretic resistance}

The resistance model follows Almeida \emph{et al.}~\cite{almeida2002hysteretic}
in the form implemented by Zhang \emph{et al.}:

\begin{align}
    R(T,\mathcal H)
    &=R_m+R_0\exp\!\left(\frac{E_a/k_B}{T}\right)g(T,\mathcal H),
    \label{eq:almeida_resistance}\\
    g_{\mathrm{major}}(T,\delta)
    &=\frac{1}{2}+\frac{1}{2}\tanh\!\left[
        \beta\left(\delta\frac{w}{2}+T_c-T\right)
    \right],
    \qquad
    \delta=\begin{cases}+1,&\text{heating},\\-1,&\text{cooling}.
    \end{cases}
    \label{eq:almeida_major}
\end{align}

Here $g$ is the semiconducting fraction, $R_m$ is the metallic resistance,
and the two values of $\delta$ define the heating and cooling branches. At a
reversal temperature $T_r$, the model chooses $T_{pr}$ to keep $g$ continuous
and evaluates the same branch at a shifted temperature:

\begin{equation}
    g(T,\mathcal H)=g_{\mathrm{major}}(T+T_p,\delta),
    \qquad
    T_p=T_{pr}P(x),
    \qquad
    x=\frac{T-T_r}{T_{pr}}.
    \label{eq:almeida_minor_loop}
\end{equation}

The shift is controlled by the proximity function

\begin{equation}
    P(x)=\frac{1}{2}\left[1-\sin(\gamma x)\right]
    \left[1+\tanh\!\left(\pi^2-2\pi x\right)\right].
    \label{eq:almeida_proximity}
\end{equation}

Thus $P(x)$ enters the resistance through the temperature shift $T_p$: it is
approximately one at the reversal and decreases as the trajectory moves away,
allowing the minor loop to approach the new major branch smoothly. The
parameter $\gamma$ controls this approach.
